\chapter{Processus ponctuels et processus de Hawkes}\label{ch:qm:point-processes-and-hawkes-processes}

Comptez les transactions sur un contrat à terme liquide pendant une matinée animée, par tranches d'une minute. Si les transactions
arrivaient comme une horloge de Poisson, la variance des comptages serait égale à leur moyenne~; elle est plusieurs
fois plus grande, et les transactions arrivent par rafales, chaque transaction rendant la suivante plus probable dans
la seconde environ. Un modèle dans lequel chaque événement relève le taux des événements ultérieurs, le
\omterm{def:qm:point-processes-and-hawkes-processes:hawkes}{processus de Hawkes}, reproduit ce phénomène, et un seul nombre le résume~: le \omterm{def:qm:point-processes-and-hawkes-processes:hawkes}{ratio de branchement}, soit le nombre moyen
de transactions que déclenche chaque transaction. Ajusté sur les contrats à terme E-mini S\&P, il a été rapporté proche de
un (Hardiman, Bercot et Bouchaud, 2013) et en hausse, passant d'environ 0.3 à plus de 0.7 en une décennie (Filimonov et
Sornette, 2012), et les mêmes auteurs ont montré avec quelle facilité un ajustement fait apparaître une excitation
absente. Ce chapitre définit les \omterm{def:qm:point-processes-and-hawkes-processes:pp}{processus ponctuels} par leurs intensités, construit le \omterm{def:qm:point-processes-and-hawkes-processes:hawkes}{processus de Hawkes}, puis
l'ajuste par maximum de vraisemblance en temps linéaire, le simule de deux façons, le vérifie par un changement d'horloge, et
mesure comment un profil intrajournalier en U trompe l'ajustement.

\section{Intensités et compensateurs}

\begin{definition}[Processus ponctuel, processus de comptage]\label{def:qm:point-processes-and-hawkes-processes:pp}
Un \emph{processus ponctuel}\index{processus ponctuel} sur $[0, \infty)$ est un ensemble aléatoire localement fini de dates
d'événements $t_1 < t_2 < \dots$. Son \emph{processus de comptage}\index{processus de comptage} est $N_t = \#\{i : t_i \le t\}$,
un \omterm{def:qm:probability-at-speed:filtration}{processus adapté}, continu à droite, à valeurs entières, non décroissant et dont les sauts valent un.
\end{definition}

\begin{definition}[Intensité conditionnelle, compensateur]\label{def:qm:point-processes-and-hawkes-processes:intensity}
L'\emph{intensité conditionnelle}\index{intensité conditionnelle} d'un \omterm{def:qm:point-processes-and-hawkes-processes:pp}{processus de comptage} est le \omterm{def:qm:probability-at-speed:filtration}{processus adapté}
$\lambda_t$ tel que $\P(N_{t+dt} - N_t = 1 \mid \mathcal F_t) = \lambda_t\,dt + o(dt)$~: le taux instantané d'un événement
connaissant le passé. Son \emph{compensateur}\index{compensateur} est $\Lambda_t = \int_0^t\lambda_s\,ds$, le processus prévisible
croissant tel que $N_t - \Lambda_t$ soit une \omterm{def:qm:ito-calculus:local}{martingale locale}.
\end{definition}

La décomposition de Doob--Meyer garantit que tout \omterm{def:qm:point-processes-and-hawkes-processes:pp}{processus de comptage} admet un \omterm{def:qm:point-processes-and-hawkes-processes:intensity}{compensateur}~; lorsqu'il est
absolument continu, sa densité est l'intensité. L'intensité est le $\lambda_t$ de la série~: le taux de Poisson
du chapitre~6, l'intensité de défaut d'un temps de défaut du livre~2, et l'intensité de Hawkes
ci-dessous sont des cas particuliers.

\begin{definition}[Processus de Poisson inhomogène, processus de Cox]\label{def:qm:point-processes-and-hawkes-processes:cox}
Un \emph{processus de Poisson inhomogène}\index{processus de Poisson inhomogène} admet une intensité déterministe
$\lambda(t)$~: les comptages sur des intervalles disjoints sont des variables de Poisson indépendantes de moyennes $\int\lambda(t)\,dt$. Un \emph{processus de
Cox}\index{processus de Cox} (\omterm{def:qm:jump-processes:poisson}{processus de Poisson} doublement stochastique) est un \omterm{def:qm:jump-processes:poisson}{processus de Poisson} dont l'intensité est
elle-même un processus aléatoire, indépendant des événements qu'il pilote.
\end{definition}

Les comptages d'un \omterm{def:qm:point-processes-and-hawkes-processes:cox}{processus de Cox} sont surdispersés~: $\Var(N_t) = \E[\Lambda_t] + \Var(\Lambda_t)$. Le livre~6 modélise les défauts
de cette façon, avec une intensité de défaut aléatoire. L'aléa d'une intensité de Cox vient de l'extérieur~; l'intensité de Hawkes
est pilotée par les propres événements du processus.

\section{Auto- et excitation mutuelle}

\begin{definition}[Processus de Hawkes, noyau d'excitation, ratio de branchement]\label{def:qm:point-processes-and-hawkes-processes:hawkes}
Un \emph{processus de Hawkes}\index{processus de Hawkes} est un \omterm{def:qm:point-processes-and-hawkes-processes:pp}{processus de comptage} d'intensité
\[
\lambda_t = \mu + \sum_{t_i < t}g(t - t_i),
\],
avec une intensité de fond $\mu > 0$ et un \emph{noyau d'excitation}\index{noyau d'excitation} $g$ positif ou nul. Le
\emph{ratio de branchement}\index{ratio de branchement} est $\lVert g\rVert_1 = \int_0^\infty g(u)\,du$. Avec le noyau
exponentiel $g(u) = \alpha e^{-\beta u}$, il vaut $\alpha/\beta$.
\end{definition}

\begin{proposition}[Stationnarité et intensité moyenne]\label{prop:qm:point-processes-and-hawkes-processes:mean}
Si $\lVert g\rVert_1 < 1$, il existe un \omterm{def:qm:point-processes-and-hawkes-processes:hawkes}{processus de Hawkes} stationnaire dont l'intensité moyenne est $\bar\lambda =
\mu/(1 - \lVert g\rVert_1)$.
\end{proposition}

\begin{proof}
En prenant l'espérance de la définition, avec $\E[dN_s] = \bar\lambda\,ds$ dans le régime stationnaire,
$\bar\lambda = \mu + \int_{-\infty}^t g(t - s)\bar\lambda\,ds = \mu + \bar\lambda\lVert g\rVert_1$. L'existence est donnée par la
construction en grappes ci-dessous, dont les grappes sont finies lorsque $\lVert g\rVert_1 < 1$.
\end{proof}

\begin{theorem}[Représentation en grappes]\label{thm:qm:point-processes-and-hawkes-processes:cluster}
Un \omterm{def:qm:point-processes-and-hawkes-processes:hawkes}{processus de Hawkes} stationnaire est un processus de grappes de Poisson~: les immigrants arrivent selon un \omterm{def:qm:jump-processes:poisson}{processus de Poisson} de taux
$\mu$~; chaque événement, immigrant ou non, a un nombre de descendants de loi de Poisson, de moyenne $\lVert g\rVert_1$, nés
après des délais indépendants de densité $g/\lVert g\rVert_1$.
\end{theorem}

\admitted

Hawkes et Oakes (1974) l'ont démontré. Chaque immigrant engendre un arbre de Galton--Watson dont la taille moyenne est
$1/(1 - \lVert g\rVert_1)$, donc une part $\lVert g\rVert_1$ des événements est constituée de descendants~: le \omterm{def:qm:point-processes-and-hawkes-processes:hawkes}{ratio de branchement} est la
fraction de l'activité qui est endogène (\cref{fig:qm:point-processes-and-hawkes-processes:cluster}). À
0.7, chaque transaction exogène en produit 3.3 au total. Les comptages le montrent~: pour le noyau exponentiel, le
rapport variance sur moyenne des comptages dans une fenêtre de longueur $\tau$ passe de un à $1/(1 - \lVert g\rVert_1)^2$ à mesure que la
fenêtre s'élargit, à partir de la densité de covariance $c(u) = \bar\lambda\frac{n(2-n)\beta}{2(1-n)}e^{-\beta(1-n)|u|}$ avec $n =
\alpha/\beta$ (\cref{fig:qm:point-processes-and-hawkes-processes:dispersion}).

\begin{omfigure}
\begin{tikzpicture}[x=1.25cm,y=0.9cm,font=\small]
\draw[->,thick,gray] (0,0) -- (11,0) node[right,black] {temps};
\foreach \x/\n in {0.8/a,4.3/b,8.2/c} {\draw[omDef,very thick] (\x,0) -- (\x,1.1); \fill[omDef] (\x,1.1) circle (2.5pt);}
\foreach \x/\p/\h in {1.4/0.8/0.8,1.7/0.8/0.8,2.3/1.7/0.55,2.6/1.4/0.55,3.1/2.6/0.35,4.9/4.3/0.8,5.8/4.9/0.55,8.5/8.2/0.8,8.9/8.2/0.8,9.4/8.9/0.55} {
  \draw[omThm,thick] (\x,0) -- (\x,\h); \fill[omThm] (\x,\h) circle (2pt);}
\draw[gray,-{Stealth}] (0.8,1.1) to[bend left=30] (1.4,0.85);
\draw[gray,-{Stealth}] (0.8,1.1) to[bend left=30] (1.7,0.85);
\draw[gray,-{Stealth}] (1.7,0.85) to[bend left=30] (2.3,0.6);
\draw[gray,-{Stealth}] (1.4,0.85) to[bend left=40] (2.6,0.6);
\draw[gray,-{Stealth}] (2.6,0.6) to[bend left=30] (3.1,0.4);
\draw[gray,-{Stealth}] (4.3,1.1) to[bend left=30] (4.9,0.85);
\draw[gray,-{Stealth}] (4.9,0.85) to[bend left=30] (5.8,0.6);
\draw[gray,-{Stealth}] (8.2,1.1) to[bend left=30] (8.5,0.85);
\draw[gray,-{Stealth}] (8.2,1.1) to[bend left=30] (8.9,0.85);
\draw[gray,-{Stealth}] (8.9,0.85) to[bend left=30] (9.4,0.6);
\node[omDef,font=\footnotesize,anchor=south] at (4.3,1.25) {immigrant (taux $\mu$)};
\node[omThm,font=\footnotesize,anchor=north] at (5.8,-0.15) {descendants~: $\text{Poisson}(\lVert g\rVert_1)$ par événement, délais $\propto g$};
\end{tikzpicture}

\omcaption{La représentation en grappes d'un \omterm{def:qm:point-processes-and-hawkes-processes:hawkes}{processus de Hawkes}. Les immigrants (hauts, en bleu) arrivent selon un \omterm{def:qm:jump-processes:poisson}{processus de
Poisson}~; chaque événement a un nombre de descendants (courts, en rouge) de loi de Poisson, de moyenne égale au \omterm{def:qm:point-processes-and-hawkes-processes:hawkes}{ratio de branchement},
à des délais tirés du noyau normalisé. Les événements de toutes les générations réunis forment le
\omterm{def:qm:point-processes-and-hawkes-processes:hawkes}{processus de Hawkes}.}\label{fig:qm:point-processes-and-hawkes-processes:cluster}
\end{omfigure}

\begin{definition}[Processus de Hawkes multivarié]\label{def:qm:point-processes-and-hawkes-processes:multi}
Un \emph{processus de Hawkes multivarié}\index{processus de Hawkes multivarié} possède les composantes $N^{(1)}, \dots,
N^{(d)}$ d'intensités $\lambda^{(i)}_t = \mu_i + \sum_j\sum_{t^{(j)}_k < t}g_{ij}(t - t^{(j)}_k)$~; la matrice de branchement
est $G = (\lVert g_{ij}\rVert_1)$, et une version stationnaire existe lorsque son rayon spectral est inférieur à un.
\end{definition}

Les achats stimulent les ventes, les transactions stimulent les révisions de cotation, l'activité d'une plateforme stimule celle d'une autre~: l'excitation mutuelle
est la façon dont sont écrits les modèles de flux d'ordres des livres~10 et~11.

\begin{example}[Achats et ventes qui s'excitent mutuellement]\label{ex:qm:point-processes-and-hawkes-processes:bivariate}
Supposons que les achats et les ventes aient chacun un niveau de base de $0.1$ par seconde et des noyaux exponentiels avec $\beta = 1$, un achat faisant
croître l'intensité d'achat de $0.2$ et l'intensité de vente de $0.4$, et symétriquement. La matrice de branchement $G = \bigl(\begin{smallmatrix}0.2 &
0.4\\0.4 & 0.2\end{smallmatrix}\bigr)$ a pour valeurs propres $0.6$ et $-0.2$, donc le processus est stable~; en prenant les espérances comme dans
la \cref{prop:qm:point-processes-and-hawkes-processes:mean}, les intensités moyennes vérifient $\bar\lambda = \mu + G\bar\lambda$, ce qui donne
$\bar\lambda = (I_2 - G)^{-1}\mu = (0.25, 0.25)$ par seconde. L'excitation croisée domine~: parmi l'activité déclenchée de chaque côté,
les deux tiers proviennent des transactions de l'autre côté.
\end{example}

\begin{omfigure}
\begin{tikzpicture}
\begin{groupplot}[group style={group size=2 by 1,horizontal sep=1.6cm},
  width=7.6cm,height=4.8cm,tick label style={font=\small},label style={font=\small},
  grid=major,grid style={gray!25},table/col sep=comma]
\nextgroupplot[xlabel={secondes},ylabel={$\lambda_t$ (par seconde)},xmin=0,xmax=60,ymin=-0.3,ymax=2.1]
\addplot[omDef,thick] table[x=s,y=lam]{figdata/methods/07-point-processes-and-hawkes-processes/intensity.csv};
\addplot[omThm,only marks,mark=|,mark size=3pt] table[x=s,y=y]{figdata/methods/07-point-processes-and-hawkes-processes/events.csv};
\addplot[gray,dashed,domain=0:60] {0.1};
\nextgroupplot[xlabel={largeur des intervalles $\tau$ (secondes)},ylabel={$\Var N(\tau)/\E N(\tau)$},xmode=log,xmin=0.2,xmax=400,ymin=0,ymax=12,
  legend style={font=\footnotesize,at={(0.03,0.97)},anchor=north west,fill=white,draw=none}]
\addplot[omDef,very thick] table[x=tau,y=theory]{figdata/methods/07-point-processes-and-hawkes-processes/dispersion.csv};
\addplot[omThm,only marks,mark=*,mark size=1.6pt] table[x=tau,y=sim]{figdata/methods/07-point-processes-and-hawkes-processes/dispersion.csv};
\addplot[gray,dashed,domain=0.2:400] {1};
\legend{théorie,simulation}
\end{groupplot}
\end{tikzpicture}

\omcaption{À gauche~: une minute d'une journée de négociation simulée avec $\mu = 0.1$, $\alpha = 0.7$ et $\beta = 1$ par seconde
(\omterm{def:qm:point-processes-and-hawkes-processes:hawkes}{ratio de branchement} 0.7)~: l'intensité saute de $\alpha$ à chaque événement (repères) et décroît vers le niveau de base
(en tirets). À droite~: rapport variance sur moyenne des comptages dans des intervalles de longueur $\tau$, issu de la densité de covariance
et dans la journée simulée, à comparer à la valeur un d'un \omterm{def:qm:jump-processes:poisson}{processus de Poisson} (en tirets), qui monte jusqu'à $1/(1 - 0.7)^2 = 11.1$. Données~: le tutoriel du chapitre,
avec graine fixée.}\label{fig:qm:point-processes-and-hawkes-processes:dispersion}
\end{omfigure}

\section{Estimation par vraisemblance}

\begin{proposition}[Log-vraisemblance d'un processus ponctuel]\label{prop:qm:point-processes-and-hawkes-processes:loglik}
Pour des événements $t_1 < \dots < t_n$ sur $[0, T]$ et une intensité paramétrique $\lambda_t(\theta)$,
\[
\ell_n(\theta) = \sum_{i=1}^n\ln\lambda_{t_i}(\theta) - \int_0^T\lambda_t(\theta)\,dt.
\]
Pour le \omterm{def:qm:point-processes-and-hawkes-processes:hawkes}{processus de Hawkes} exponentiel, avec $A_i = \sum_{j<i}e^{-\beta(t_i - t_j)}$,
\[
\ell_n = \sum_i\ln(\mu + \alpha A_i) - \mu T - \frac\alpha\beta\sum_i\bigl(1 - e^{-\beta(T - t_i)}\bigr), \qquad A_i =
e^{-\beta(t_i - t_{i-1})}(1 + A_{i-1}).
\]
\end{proposition}

\begin{proof}
La densité des événements est le produit, sur les intervalles, de la probabilité qu'il n'y ait aucun événement, $e^{-\int\lambda}$,
multipliée par l'intensité en chaque événement, $\lambda_{t_i}$~: la vraisemblance d'une suite de temps d'attente exponentiels
à intensité variable dans le temps. La récurrence découle de $e^{-\beta(t_i - t_j)} = e^{-\beta(t_i - t_{i-1})}e^{-\beta(t_{i-1} - t_j)}$.
\end{proof}

La récurrence rend le calcul de la vraisemblance linéaire en nombre d'événements, au lieu de quadratique~; c'est le
noyau de l'estimateur du projet en cours, en Python, C++20 et Rust (le
second listing du tutoriel). Le maximum est trouvé numériquement~: sur une
journée simulée de 8\,006 événements avec $n = 0.7$, l'estimation donne $(\hat\mu, \hat\alpha, \hat\beta) = (0.103, 0.696, 0.995)$,
soit un \omterm{def:qm:point-processes-and-hawkes-processes:hawkes}{ratio de branchement} de 0.700~; sur quarante journées simulées, les estimations ont une moyenne de 0.700 et un écart-type
de 0.011. Le chapitre~11 dérive les erreurs types qui accompagnent le maximum de vraisemblance.

\section{Simulation par amincissement et par branchement}

\begin{definition}[Amincissement]\label{def:qm:point-processes-and-hawkes-processes:thinning}
L'\emph{amincissement}\index{amincissement} simule un \omterm{def:qm:point-processes-and-hawkes-processes:pp}{processus ponctuel} d'intensité $\lambda_t \le M_t$, où $M$ est une borne connue
à l'avance, en proposant des événements au taux $M$ et en acceptant une proposition en $t$ avec la probabilité $\lambda_t/M_t$.
\end{definition}

\begin{method}[Amincissement d'Ogata pour un processus de Hawkes]\label{met:qm:point-processes-and-hawkes-processes:ogata}
Entre les événements, l'intensité de Hawkes exponentielle décroît, de sorte que sa valeur juste après le dernier événement la borne
jusqu'au suivant. Tirer un temps d'attente à ce taux, faire décroître l'excitation sur ce délai, accepter avec probabilité
$\lambda/\text{borne}$, et en cas d'acceptation ajouter $\alpha$ à l'excitation. La simulation est exacte, sans pas de temps (le
premier listing du tutoriel). Autrement, simuler les grappes du
\cref{thm:qm:point-processes-and-hawkes-processes:cluster} génération par génération~: la part des descendants
parmi les événements est alors observable, 0.699 sur une journée simulée contre 0.7.
\end{method}

\section{Diagnostics par changement de temps}

\begin{theorem}[Changement d'échelle temporelle]\label{thm:qm:point-processes-and-hawkes-processes:rescaling}
Si $N$ admet un \omterm{def:qm:point-processes-and-hawkes-processes:intensity}{compensateur} continu $\Lambda$ avec $\Lambda_\infty = \infty$, les temps rééchelonnés $\Lambda(t_1) < \Lambda(t_2) < \dots$
forment un \omterm{def:qm:jump-processes:poisson}{processus de Poisson} de taux unité~: les écarts $\Lambda(t_i) - \Lambda(t_{i-1})$ sont indépendants et de loi $\mathrm{Exp}(1)$.
\end{theorem}

\admitted

En faisant défiler l'horloge à la vitesse de l'intensité du modèle, tout \omterm{def:qm:point-processes-and-hawkes-processes:pp}{processus ponctuel} bien spécifié devient un
\omterm{def:qm:jump-processes:poisson}{processus de Poisson} standard, de sorte que le \omterm{def:qm:point-processes-and-hawkes-processes:intensity}{compensateur} ajusté fournit des résidus à tester~: leur moyenne doit valoir un, leur
variance doit valoir un, leurs quantiles ceux de la loi exponentielle (Brown et al., 2002). Pour le modèle de Hawkes ajusté
les résidus ont une moyenne de 1.000 et une variance de 0.98~; pour un modèle de Poisson de même taux moyen, la variance vaut
4.63, et les quantiles supérieurs sont bien trop grands (la \cref{fig:qm:point-processes-and-hawkes-processes:qq}).

\begin{omfigure}
\begin{tikzpicture}
\begin{axis}[width=8.5cm,height=6cm,
  xlabel={quantile de $\mathrm{Exp}(1)$},ylabel={quantile des résidus},
  tick label style={font=\small},label style={font=\small},
  xmin=0,xmax=10,ymin=0,ymax=28,grid=major,grid style={gray!25},
  legend style={font=\footnotesize,at={(0.03,0.97)},anchor=north west,fill=white,draw=none},
  table/col sep=comma]
\addplot[omThm,thick,mark=*,mark size=1.1pt] table[x=q,y=poisson]{figdata/methods/07-point-processes-and-hawkes-processes/qq.csv};
\addplot[omDef,thick,mark=*,mark size=1.1pt] table[x=q,y=hawkes]{figdata/methods/07-point-processes-and-hawkes-processes/qq.csv};
\addplot[gray,dashed,domain=0:10] {x};
\legend{ajustement de Poisson,ajustement de Hawkes,droite à 45 degrés}
\end{axis}
\end{tikzpicture}

\omcaption{Quantiles des résidus rééchelonnés en temps d'une journée simulée, comparés à ceux d'une
exponentielle unitaire. Le modèle de Hawkes ajusté se situe sur la diagonale~; un modèle de Poisson de même taux moyen manque les
rafales, et ses plus grands intervalles sont presque trois fois trop longs. Données~: le tutoriel du chapitre,
avec graine fixée.}\label{fig:qm:point-processes-and-hawkes-processes:qq}
\end{omfigure}

\section{Tutoriel~: la bande auto-excitée}\label{tut:qm:point-processes-and-hawkes-processes:tape}

\begin{tutorial}
\textbf{Objectif.} Simuler une journée de transactions auto-excitées, l'ajuster par maximum de vraisemblance, vérifier l'ajustement
par rééchelonnement en temps, et ajuster une journée sans excitation mais avec une configuration infrajournalière. \textbf{État final~:}
\cref{fig:qm:point-processes-and-hawkes-processes:dispersion,fig:qm:point-processes-and-hawkes-processes:qq}
et les \omterm{def:qm:point-processes-and-hawkes-processes:hawkes}{ratios de branchement} du problème du week-end.
\begin{tutsteps}
\item \textbf{Simuler} par l'\omterm{def:qm:point-processes-and-hawkes-processes:thinning}{amincissement} d'Ogata.
\omcode{code/firm/hawkes/firm_hawkes.py}{24}{40}{Amincissement d'Ogata pour le processus de Hawkes exponentiel.}\label{lst:qm:point-processes-and-hawkes-processes:thinning}
\item \textbf{Ajuster}~: la log-vraisemblance en temps linéaire, maximisée par une recherche de Nelder--Mead sur $(\ln\mu,
\operatorname{logit}n, \ln\beta)$ de sorte que le \omterm{def:qm:point-processes-and-hawkes-processes:hawkes}{ratio de branchement} reste inférieur à un.
\omcode{code/firm/hawkes/firm_hawkes.py}{98}{111}{La log-vraisemblance avec la récurrence pour $A_i$.}\label{lst:qm:point-processes-and-hawkes-processes:loglik}
\item \textbf{Exécuter} \texttt{problem()} et \texttt{fig\_hawkes.py}~; les jumeaux C++20 et Rust dans
\texttt{code/firm/hawkes/} reproduisent la log-vraisemblance Python à $10^{-12}$ sur une liste d'événements fixée.
\end{tutsteps}
\textbf{Pour aller plus loin.} Ajuster un noyau à deux exponentielles et voir si les parties rapide et lente se séparent~;
simuler un processus bivarié d'achats et de ventes qui s'excitent mutuellement et retrouver la matrice de branchement.
\end{tutorial}

\section{Construction~: la boîte à outils de Hawkes}\label{bld:qm:point-processes-and-hawkes-processes:hawkes}

\begin{build}
\bfield{Rôle} Modéliser l'arrivée des transactions, des ordres et des annulations sur les marchés de la firme miniature~:
simuler un flux réaliste et par rafales pour le simulateur de bourse, et estimer l'excitation
à partir des flux enregistrés.
\bfield{Interface} Python~: \texttt{simulate\_thinning}, \texttt{simulate\_branching},
\texttt{simulate\_multivariate}, \texttt{intensity}, \texttt{loglik}, \texttt{fit}, \texttt{compensator},
\texttt{residuals}, \texttt{nelder\_mead}. C++20 et Rust~: \texttt{Params}, \texttt{simulate\_thinning},
\texttt{loglik}, \texttt{compensator}.
\bfield{Règles} Simulation exacte (sans pas de temps)~; vraisemblance en $O(n)$~; paramètres recherchés sur une échelle qui
conserve $\mu, \beta > 0$ et $0 \le n < 1$~; les trois langages concordent à $10^{-12}$ sur la liste de référence.
\bfield{Tests d'acceptation} \texttt{code/firm/hawkes/\{tests,cpp,rust\}}~: les effectifs moyens égalent $\mu T/(1 - n)$~;
le maximum de vraisemblance retrouve les paramètres~; les résidus du vrai modèle sont des exponentielles unitaires~; les simulations par branchement et
par \omterm{def:qm:point-processes-and-hawkes-processes:thinning}{amincissement} ont la même loi de comptage.
\bfield{Pour aller plus loin} Maximum de vraisemblance multivarié~; noyaux en loi de puissance~; une référence $\mu(t)$ variable dans le temps, estimée
conjointement avec le noyau (la correction du problème du week-end).
\end{build}

\begin{omsources}
\item A.~G.~Hawkes, «~Spectra of some self-exciting and mutually exciting point processes~»,
\emph{Biometrika} 58, 1971.
\item A.~G.~Hawkes et D.~Oakes, «~A cluster process representation of a self-exciting process~»,
\emph{Journal of Applied Probability} 11, 1974.
\item Y.~Ogata, «~On Lewis' simulation method for point processes~», \emph{IEEE Transactions on
Information Theory} 27, 1981.
\item E.~N.~Brown, R.~Barbieri, V.~Ventura, R.~E.~Kass et L.~M.~Frank, «~The time-rescaling theorem and
its application to neural spike train data analysis~», \emph{Neural Computation} 14, 2002.
\item V.~Filimonov et D.~Sornette, «~Quantifying reflexivity in financial markets~», \emph{Physical Review E}
85, 2012~; «~Apparent criticality and calibration issues in the Hawkes self-excited point process model~»,
\emph{Quantitative Finance} 15, 2015.
\item S.~J.~Hardiman, N.~Bercot et J.-P.~Bouchaud, «~Critical reflexivity in financial markets: a Hawkes
process analysis~», \emph{European Physical Journal B} 86, 2013.
\end{omsources}

\section{Exercices}

\begin{exercise}[$\star$]\label{exo:qm:point-processes-and-hawkes-processes:1}
Un \omterm{def:qm:point-processes-and-hawkes-processes:hawkes}{processus de Hawkes} a $\mu = 0.1$ par seconde, $\alpha = 0.7$ et $\beta = 1$ par seconde. Quels sont son \omterm{def:qm:point-processes-and-hawkes-processes:hawkes}{ratio de branchement}, son intensité moyenne
et le nombre d'événements attendu sur une journée de négociation de 23\,400 secondes~?
\end{exercise}

\begin{exercise}[$\star$]\label{exo:qm:point-processes-and-hawkes-processes:2}
Dans le même processus, combien d'événements un immigrant engendre-t-il en moyenne, et quelle part des événements sont
des immigrants~?
\end{exercise}

\begin{exercise}[$\star$]\label{exo:qm:point-processes-and-hawkes-processes:3}
Un \omterm{def:qm:point-processes-and-hawkes-processes:cox}{processus de Poisson inhomogène} a une intensité $\lambda(t) = a + bt$. Écrire son \omterm{def:qm:point-processes-and-hawkes-processes:intensity}{compensateur}, et la probabilité
d'absence d'événement sur $[0, T]$.
\end{exercise}

\begin{exercise}[$\star\star$]\label{exo:qm:point-processes-and-hawkes-processes:4}
Au bout de combien de temps après un événement sa contribution à l'intensité met-elle à se réduire de moitié, avec $\beta = 1$ par seconde~? Quelle est
l'ampleur du saut de l'intensité à un événement par rapport à la référence de base~?
\end{exercise}

\begin{exercise}[$\star\star$]\label{exo:qm:point-processes-and-hawkes-processes:5}
Calculer le rapport variance sur moyenne des comptages dans des intervalles de 1 et de 60 secondes, ainsi que sa limite, pour le processus de
l'\cref{exo:qm:point-processes-and-hawkes-processes:1}.
\end{exercise}

\begin{exercise}[$\star\star$]\label{exo:qm:point-processes-and-hawkes-processes:6}
Chaque jour, le taux d'arrivée vaut 1 ou 3 par seconde avec la même probabilité, et à l'intérieur d'une journée les arrivées suivent un \omterm{def:qm:jump-processes:poisson}{processus de Poisson}.
Quelles sont la moyenne et la variance du nombre de comptages en une seconde~? S'agit-il d'un \omterm{def:qm:point-processes-and-hawkes-processes:hawkes}{processus de Hawkes}~?
\end{exercise}

\begin{exercise}[$\star\star\star$]\label{exo:qm:point-processes-and-hawkes-processes:7}
\emph{Programmation.} Ajuster le modèle de Hawkes exponentiel à quarante journées simulées du processus de
l'\cref{exo:qm:point-processes-and-hawkes-processes:1} et indiquer la moyenne et l'écart type
du \omterm{def:qm:point-processes-and-hawkes-processes:hawkes}{ratio de branchement} estimé.
\end{exercise}

\begin{exercise}[$\star\star\star$]\label{exo:qm:point-processes-and-hawkes-processes:8}
\emph{Repérer l'erreur.} « Nous avons ajusté un modèle de Hawkes à une journée complète de transactions et trouvé un \omterm{def:qm:point-processes-and-hawkes-processes:hawkes}{ratio de branchement} de 0.87~:
le marché est proche de la criticité. »
\end{exercise}

\section{Problème~: la bande auto-excitée}

\begin{problem}[{Problème du week-end --- une excitation présente et une excitation absente}]\label{pb:qm:point-processes-and-hawkes-processes:1}
Les transactions sur un contrat arrivent pendant une journée de 23\,400 secondes comme un \omterm{def:qm:point-processes-and-hawkes-processes:hawkes}{processus de Hawkes} exponentiel avec $\mu = 0.1$, $\alpha =
0.7$ et $\beta = 1$ par seconde. Un second jour ne présente aucune excitation~: les transactions forment un \omterm{def:qm:point-processes-and-hawkes-processes:cox}{processus de Poisson inhomogène}
dont le taux est trois fois plus élevé à l'ouverture et à la clôture qu'à midi, $\lambda(t) \propto 1 + 8(t/T -
\tfrac12)^2$, avec le même nombre moyen de transactions.

\medskip\noindent\textbf{Part I --- The model.}
\begin{enumerate}
\item Quel est le \omterm{def:qm:point-processes-and-hawkes-processes:hawkes}{ratio de branchement}~?
\item Quels sont l'intensité moyenne et le nombre moyen de transactions de la journée~?
\item Quelle est la taille moyenne d'un amas, et quelle est la part des immigrants~?
\item Quelle est la \omterm{def:qm:stochastic-differential-equations:ou}{demi-vie} de l'excitation d'un événement~?
\item Quels sont les rapports variance sur moyenne des comptages à la seconde et à la minute~?
\end{enumerate}

\medskip\noindent\textbf{Partie II --- Estimation.}
\begin{enumerate}[resume]
\item Combien de transactions compte la journée simulée~?
\item Qu'estime le maximum de vraisemblance~?
\item Sur quarante journées simulées, quelles sont la moyenne et l'écart type du \omterm{def:qm:point-processes-and-hawkes-processes:hawkes}{ratio de branchement} estimé~?
\item Quelles moyenne et variance ont les résidus rééchelonnés dans le temps, pour l'ajustement de Hawkes et pour un ajustement de Poisson~?
\item Dans une simulation par branchement de la journée, quelle part des événements sont des enfants~?
\end{enumerate}

\medskip\noindent\textbf{Partie III --- La journée saisonnière.}
\begin{enumerate}[resume]
\item Quel est le taux à l'ouverture par rapport à midi, et le taux moyen~?
\item Combien de transactions compte la journée saisonnière~?
\item Quels \omterm{def:qm:point-processes-and-hawkes-processes:hawkes}{ratio de branchement} et taux de décroissance l'ajustement de Hawkes trouve-t-il~?
\item Sur vingt journées saisonnières, puis sur vingt journées d'un \omterm{def:qm:jump-processes:poisson}{processus de Poisson} à taux constant~?
\item Pourquoi l'ajustement trouve-t-il de l'excitation là où il n'y en a pas~?
\end{enumerate}

\medskip\noindent\textbf{Partie IV --- Jugement.}
\begin{enumerate}[resume]
\item Comment supprimeriez-vous le biais~?
\item Que suggèrent les estimations publiées contradictoires pour le E-mini~?
\item Comment un bureau de négociation pourrait-il utiliser un modèle de Hawkes bien ajusté~?
\item Énoncez le \emph{résultat nommé}~: le \omterm{def:qm:point-processes-and-hawkes-processes:hawkes}{ratio de branchement} retrouvé, et le ratio fallacieux.
\item En une phrase~: quand un \omterm{def:qm:point-processes-and-hawkes-processes:hawkes}{ratio de branchement} élevé est-il la preuve d'une auto-excitation~?
\end{enumerate}
\end{problem}

\section{Questions d'entretien}

\begin{interviewq}[$\star$ \iqroles{chercheur, trader}]\label{iq:qm:point-processes-and-hawkes-processes:1}
Qu'est-ce qu'un \omterm{def:qm:point-processes-and-hawkes-processes:hawkes}{processus de Hawkes}, et que mesure son \omterm{def:qm:point-processes-and-hawkes-processes:hawkes}{ratio de branchement}~?
\end{interviewq}

\begin{interviewq}[$\star\star$ \iqroles{chercheur}]\label{iq:qm:point-processes-and-hawkes-processes:2}
Montrer que l'intensité moyenne d'un \omterm{def:qm:point-processes-and-hawkes-processes:hawkes}{processus de Hawkes} stationnaire vaut $\mu/(1 - n)$. Que se passe-t-il lorsque $n \ge 1$~?
\end{interviewq}

\begin{interviewq}[$\star\star$ \iqroles{chercheur, développeur}]\label{iq:qm:point-processes-and-hawkes-processes:3}
Comment simule-t-on un \omterm{def:qm:point-processes-and-hawkes-processes:hawkes}{processus de Hawkes} de façon exacte~?
\end{interviewq}

\begin{interviewq}[$\star\star$ \iqroles{développeur}]\label{iq:qm:point-processes-and-hawkes-processes:4}
La log-vraisemblance de Hawkes semble quadratique dans le nombre d'événements. Comment l'évaluer en temps linéaire pour un
noyau exponentiel~?
\end{interviewq}

\begin{interviewq}[$\star\star$ \iqroles{trader, chercheur}]\label{iq:qm:point-processes-and-hawkes-processes:5}
Comment un teneur de marché utiliserait-il une estimation de l'auto-excitation des transactions~?
\end{interviewq}

\begin{interviewq}[$\star\star\star$ \iqroles{chercheur, ingénieur ML}]\label{iq:qm:point-processes-and-hawkes-processes:6}
Comment testeriez-vous si un modèle de \omterm{def:qm:point-processes-and-hawkes-processes:pp}{processus ponctuel} ajusté est adéquat~?
\end{interviewq}
